\subsection{Sprint 4}

Para esta sprint estavam previstas 7 User Stories, a maioria tratando de autenticação, login, usuários, JWT e observabilidade com tasks relativas a DevOps. Da minha parte fiquei responsável por garantir que o respositório do Backend estivesse perfeito para entrega final e fazer algumas tasks de DevOps. Além disso fui o único integrante responsável por garantir o deploy do backend para a entrega final, que apesar de parecer uma tarefa simples, traz consigo muitos bugfixes de infra e ci/cd, principalmente num cenário de poucas sprints como a Ages que a pipeline foi testada apenas 2 vezes antes da entrega final

O time conseguiu concluir todas as tasks com eficiência desde a primeira até a última sprint do projeto, ressaltando a dedicação e proatividade dos membros. Da minha parte, consegui um bom tempo livre para realizar tarefas, principalmente devido ao fato de eu ter tido uma lesão que me afastou das aulas por 3 semanas, proporcionando mais tempo em casa para fazer tasks. Consegui aprovar todos MRs relativos ao Backend, fazer tasks de Infra, como Sentry, Prometheus, Dependabot e outras, além de ajudar com bugfixes para a entrega final. Por fim, fiz o deploy um dia antes da entrega e fico feliz em dizer que separei um bom tempo do dia para corrigir algum possível erro de deploy, pois encontrei alguns e graças a espera do imprevisto, consegui fazer com calma o deploy corrigindo problemas de espaço e dependências no container da AWS.

Como falei antes, esta sprint foi para mim, diferente de todas as outras, tendo em vista que com a lesão e o afastamento das aulas, tive que continuar me comunicando com o time mesmo estando de casa. Creio que isto teria sido um problema muito maior se o time não estivesse tão entrosado, talvez até em uma Sprint 0 ou 1, teria tido problemas nesse fator, no entanto, um pouco por sorte e principalmente por competência do time, de criar uma ambiente tão amigável, posso dizer que a comunicação, mesmo de casa, não representou um problema.

Essa Ages me ensinou a importância de uma gestão de Ages 4, percebi que houveram as tasks mais mastigadas de todas as Ages, deixando bem detalhado como o colega deveria fazer a implementação, facilitando a vida de todos que pegavam uma task para fazer e evitando que houvessem ambiguidades.

Por ser a última sprint da Ages, o foco dos próximos passos já é na próxima, que por sinal, começa em cerca de 1 mês, acho que tenho muito o que tirar de inspiração de meus Ages 4 para minha atuação da próxima Ages. Além disso uma revisão da matéria de Gerenciamento de Projeto de Software, também servirá para refrescar os conteúdos mais importantes que tenho que levar em conta na Ages 4.
